\subsection{Contrast 1: RPE}

\begin{table}[htbp]
    \centering
    \begin{threeparttable}
    \caption{Clusters that correlate with RPE, cluster-forming threshold T = 8.}
    \label{tab:con1_pos}%
    \begin{tabularx}{\textwidth}{l *{6}{R}}
        \toprule
        \textbf{Cluster Label} & \textbf{Voxels} & \textbf{MAX T} & \textbf{MAX X (mm)} & \textbf{MAX Y (mm}) & \textbf{MAX Z (mm)} & \textbf{P FWE-CORR.} \\
        \midrule
        Caudate, Putamen R & 391 & 19.29 & 18 & 14 & -12 & 0.0002 \\ % cluster 0
Caudate, Putamen L & 273 & 16.22 & -16 & 14 & -12 & 0.0002 \\ % cluster 2
Angular Gyrus L & 66 & 11.78 & -46 & -66 & 44 & 0.0002 \\ % cluster 7
Visual cortex (Mid, Lat Occipital) L & 139 & 11.78 & -46 & -82 & 2 & 0.0002 \\ % cluster 4
vmPFC & 213 & 11.05 & 8 & 54 & -8 & 0.0002 \\ % cluster 3
Visual Cortex (Lingual, Fusiform) L & 61 & 10.80 & -22 & -88 & -12 & 0.0002 \\ % cluster 8
Cerebellum R & 111 & 10.67 & 42 & -70 & -40 & 0.0002 \\ % cluster 5
        \bottomrule
    \end{tabularx}
    \end{threeparttable}
\end{table}%